% !TeX root = ../../thesis.tex
\chapter{Scene Dynamification}
\label{ch:8_Scene_Dynamification}

% ─────────────────────────────────────────────────────────────────────────────
\section{Introduction}
\label{sec:8_intro}

% Paragraph purpose: situate the chapter as the final pipeline stage and frame it by building-element class.
This chapter addresses the final stage of the pipeline, in which the static reconstructions produced by the preceding chapters are turned into a single interactive \gls{DRM}. The completed environment of Chapter~\ref{ch:6_Scene_Completion} and the completed objects of Chapter~\ref{ch:7_Object_Completion} are geometrically accurate but rigid: their walls are closed, their openings are baked into flat texture, and their objects cannot be resized without distorting their structural proportions~\cite{vermandere_adaptive_2026}. Dynamification removes this rigidity across the three classes of building element introduced in Chapter~\ref{ch:background}. The primary elements, the structural shell of walls, floor and ceiling, are made adjustable so the room envelope can be reshaped. The secondary elements, the windows and doors that the scene completion left hidden inside the inpainted wall textures, are recovered as separable, movable components. The tertiary elements, the furniture and loose objects, are converted into part-aware, selectively scalable entities. The result is a scene in which every element can be manipulated after reconstruction rather than a fixed reconstruction of a single moment (Figure~\ref{fig:8_intro}).

\iffalse
\begin{figure}[hbht]
 \centering
 \includegraphics[width=1\textwidth]{Scene_dynamification_intro}
 \caption[Scene dynamification introduction]{The static reconstructed scene (left) is converted into a dynamic representation (right) in which the structural shell is adjustable, openings are separable, and objects support part-aware scaling.}
 \label{fig:8_intro}
\end{figure}
\fi

% Paragraph purpose: state scope, what is reused and what is newly created.
The scene-level dynamification reuses the per-plane reconstructed and inpainted meshes of Chapter~\ref{ch:6_Scene_Completion}, and the object-level dynamification reuses the completion-network output of Chapter~\ref{ch:7_Object_Completion}. Beyond recombining these into one coherent dynamic scene, the chapter introduces two pieces of new geometry and texture. At the scene level, windows and doors are separated from the walls and the surface behind them is reconstructed and re-inpainted, so the wall stays complete when an opening is moved. At the object level, the static completion output is converted into the \gls{CSDF} representation on which the deformation operates. These steps are detailed in the methodology.


% ─────────────────────────────────────────────────────────────────────────────
\section{Related Work}
\label{sec:8_related}

This section reviews the four areas specific to dynamification: window and door detection, part segmentation, surface-based deformation, and implicit-field deformation. The representations and completion methods underlying the chapter were covered in Chapter~\ref{ch:background} and Chapter~\ref{ch:7_Object_Completion} and are not repeated here.


% ─────────────────────────────────────────────────────────────────────────────
\subsection{Window and Door Detection}
\label{sec:8_rw_detection}

Recovering the openings hidden in the wall textures requires detecting windows and doors in an image without a fixed, pre-trained class set, since the texture content varies widely across rooms and datasets. Open-vocabulary detectors address this by conditioning detection on a free-text prompt. Grounding DINO~\cite{liu_grounding-dino_2023} couples a transformer detector with grounded language pre-training, so that arbitrary phrases such as \emph{window} or \emph{door} can be localised in an image without category-specific training. To turn a detected box into a precise region outline, the Segment Anything Model (SAM)~\cite{kirillov_segment_2023} produces a tight mask from a box or point prompt and generalises across domains without fine-tuning. Combining a prompt-driven detector with a promptable segmenter gives the class-agnostic, mask-level localisation needed to cut openings out of the reconstructed walls, which is the combination adopted in this chapter.


% ─────────────────────────────────────────────────────────────────────────────
\subsection{Part Segmentation}
\label{sec:8_rw_segmentation}

To deform an object part by part, it must first be divided into parts. Learning-based shape segmentation methods predict a semantic or instance label per point. CSN~\cite{loizou_cross-shape_2023} introduces a cross-shape attention mechanism that lets the point features of one shape interact with those of others, improving the consistency of the predicted segmentation across a collection. Mid-Net~\cite{wang_unsupervised_2020} learns a generic shape encoding from unlabelled point clouds by jointly encoding shape and point features, and FG-Net~\cite{liu_fg-net_2020} processes large-scale point clouds directly without voxelisation through correlated feature mining and deformable convolutions. These methods target the semantic decomposition of objects, which is difficult to generalise to the arbitrary, partially reconstructed objects that arrive from Chapter~\ref{ch:7_Object_Completion}, where class labels are unreliable and the part vocabulary is open-ended.

Approximate convex decomposition (ACD) sidesteps the semantic question by splitting a shape into nearly convex components on purely geometric grounds, which is sufficient for defining deformation regions. Classical methods such as HACD~\cite{mamou_simple_2009} cluster mesh faces hierarchically using a concavity-driven merge heuristic, which is robust but tends to produce redundant and overlapping parts. Learning-based approaches such as CvxNet~\cite{deng_cvxnet_2020} represent a shape as a union of learned convex primitives, yielding compact decompositions but generalising poorly outside the training distribution. CoACD~\cite{wei_approximate_2022} cuts the triangle mesh directly with planes, scores cuts with a collision-aware concavity metric that is sensitive to interior geometry, and searches cut sequences through tree search rather than greedy splitting, producing intersection-free convex parts with fewer components. Although the resulting parts do not necessarily coincide with semantic components, their granularity and class-agnostic behaviour make them well suited to defining scaling regions on arbitrary reconstructed objects, which is why CoACD is adopted in this chapter.


% ─────────────────────────────────────────────────────────────────────────────
\subsection{Surface-Based Deformation}
\label{sec:8_rw_surface}

A second family of methods deforms the explicit surface directly. Nine-slicing~\cite{w3_css-backgrounds-3border-image-slice_2024} divides a 2D image into a grid so that corner regions stay fixed while edge and centre regions stretch, and its 3D extension, 27-slicing~\cite{deftly_27_2021}, applies the same idea to a bounding box so that an object can be scaled without distorting its corners. Both rely on manually defined zones and impose no constraint on the deformed region, which restricts them to regular, box-like shapes. This makes them a natural fit for the rectangular framing of windows and doors, where the frame should keep its thickness while the panels stretch, and 27-slicing is used for exactly that purpose in this chapter. KeypointDeformer~\cite{jakab_keypointdeformer_2021} replaces the manual zones with automatically detected keypoints that drive a cage deformation, producing more natural transformations for moderate edits but degrading under large deformations. Skeleton-driven procedural methods such as that of Getto et al.~\cite{getto_automatic_2020} detect object components from a skeleton to bound the deformation regions, improving results on more complex shapes, but they share the limitation that a stretched part is simply rescaled to fit the target size with no notion of part-size consistency. Proc-GS~\cite{li_proc-gs_2024} instead handles large deformations by treating buildings as assemblies of modular components that can be repeated rather than stretched, but it depends on a predefined library of parts. The repetition-based strategy used in this chapter borrows the idea of repeating modular components from Proc-GS while removing the dependence on a curated part library, deriving the repeatable units directly from the convex decomposition of the object instead~\cite{vermandere_adaptive_2026}.


% ─────────────────────────────────────────────────────────────────────────────
\subsection{Implicit-Field Deformation}
\label{sec:8_rw_implicit}

Because the objects of Chapter~\ref{ch:7_Object_Completion} are represented as \glspl{SDF}, the deformation must ultimately act on the implicit field. DIF-Net~\cite{deng_deformed_2021} represents a category of shapes with a shared template field together with per-instance deformation and correction fields, enabling non-destructive editing within a learned shape space. SALAD~\cite{koo_salad_2023} uses a part-level latent diffusion model for generating and editing shapes, DualSDF~\cite{hao_dualsdf_2020} exposes a coarse and a fine level of the same \gls{SDF} for semantic manipulation, and HybridSDF~\cite{vasu_hybridsdf_2022} combines implicit shapes with parametric primitives to balance flexibility against structural coherence. For part-level editing, SPAGHETTI~\cite{hertz_spaghetti_2022} supports affine transformations of implicit parts such as rotation and translation with smooth blending, and SENS~\cite{binninger_sens_2024} extends it with sketch-based control, but free scaling at the sub-part level, in particular large scaling that should preserve rather than stretch repeated structure, remains difficult and typically depends on category-specific priors. The method in this chapter targets exactly this gap: it deforms the \gls{CSDF} field of an arbitrary completed object through direct interpolation and repetition, without requiring a learned shape space or a predefined part database~\cite{vermandere_adaptive_2026}.


% ─────────────────────────────────────────────────────────────────────────────
\section{Methodology}
\label{sec:8_methodology}

The dynamification is organised into a scene-level stage and an object-level stage (Figure~\ref{fig:8_method_overview}). The scene-level stage (Section~\ref{sec:8_scene}) reopens the closed structural shell by detecting windows and doors, separating them as dynamic elements and reconstructing the surface behind them. The object-level stage (Section~\ref{sec:8_object}) converts each completed object into a \gls{CSDF}, decomposes it into convex parts, and scales user-defined regions of the field. The recombination of both into a single interactive scene is described in Section~\ref{sec:8_recombination}.

\iffalse
\begin{figure}[hbht]
 \centering
 \includegraphics[width=1\textwidth]{Scene_dynamification_method}
 \caption[Scene dynamification pipeline]{Overview of the dynamification pipeline. The completed scene (top) is reopened at its detected windows and doors, while the completed objects (bottom) are decomposed and selectively scaled, before both are recombined into a single interactive scene.}
 \label{fig:8_method_overview}
\end{figure}
\fi


% ─────────────────────────────────────────────────────────────────────────────
\subsection{Scene Dynamification}
\label{sec:8_scene}

The scene completion of Chapter~\ref{ch:6_Scene_Completion} deliberately reconstructs each room as a set of closed, solid walls, so that openings such as windows and doors appear only as flat regions painted into the wall texture by the inpainting step. Scene dynamification recovers these openings as independent, manipulable geometry. Each opening is detected in the wall texture, cut out as a separate surface, turned into a dynamic element with its own frame and panels, and the wall surface behind it is reconstructed so that the structural shell stays complete.


% ── 8.3.1.1 ──────────────────────────────────────────────────────────────────
\subsubsection{Opening Detection}
\label{sec:8_opening_detection}

Each wall fragment $M_k$ produced in Chapter~\ref{ch:6_Scene_Completion} carries a flattened texture in its own \gls{UV} space, which already presents the surface upright and free of perspective distortion~\cite{vermandere_semantic_2024}. This is the ideal input for a 2D detector, because each window and door appears as a flat, front-facing rectangle rather than a tilted and stretched shape seen from an angle. The base-colour texture of each fragment is extracted and passed to the Grounding DINO~\cite{liu_grounding-dino_2023} open-vocabulary detector, prompted with the phrases \emph{window} and \emph{door}, which returns a bounding box and class label per opening. Each box is then refined into a tight binary mask by SAM~\cite{kirillov_segment_2023}, prompted with that box, and the largest external contour of the mask is traced and simplified into a pixel-space polygon $\{(x_l, y_l)\}$ describing the opening outline (Figure~\ref{fig:8_opening_detection}).

\iffalse
\begin{figure}[hbht]
 \centering
 \includegraphics[width=1\textwidth]{Scene_dynamification_opening_detection}
 \caption[Window and door detection on wall textures]{The flattened wall texture (left) is processed by Grounding DINO~\cite{liu_grounding-dino_2023} and SAM~\cite{kirillov_segment_2023} to produce a polygonal outline of each window and door (right).}
 \label{fig:8_opening_detection}
\end{figure}
\fi

Operating per fragment, in its canonical \gls{UV} orientation, keeps the detection robust: each wall is seen in isolation and upright, so the detector is not confused by adjacent surfaces or by the curvature of the original panorama. A small padding is applied to each detection so that window and door frames, which sit at the very edge of the detected box, are included in the subsequent cutout.


% ── 8.3.1.2 ──────────────────────────────────────────────────────────────────
\subsubsection{Opening Cutout}
\label{sec:8_opening_cutout}

Each detected polygon is lifted from texture space back onto the wall mesh to extract the opening as a separate surface. The pixel polygon is first converted to \gls{UV} coordinates by normalising each vertex by the texture dimensions, $\mathbf{c}_l = (x_l / W,\, y_l / H)$, with the vertical axis flipped to match the glTF texture convention where required. The polygon is optionally expanded by a fixed margin in texture space so that the cut includes the full frame.

The wall fragment is then clipped against this \gls{UV} polygon triangle by triangle. For each mesh triangle, its three \gls{UV} coordinates define a triangle in texture space; this triangle is intersected with the opening polygon, and the intersection is re-triangulated. Each new vertex, given by its \gls{UV} position, is mapped back to 3D using the barycentric coordinates of that position within the original triangle, so that the extracted surface still samples the original wall texture and lies exactly on the wall plane. For an interior point with barycentric weights $(b_0, b_1, b_2)$ relative to the original triangle vertices $\mathbf{v}_0, \mathbf{v}_1, \mathbf{v}_2$, the 3D position is

\[
\mathbf{x} = b_0\,\mathbf{v}_0 + b_1\,\mathbf{v}_1 + b_2\,\mathbf{v}_2 .
\]

The union of all clipped triangles forms a new textured mesh containing only the opening region, separated from the wall but geometrically and texturally consistent with it (Figure~\ref{fig:8_opening_cutout}). All geometric detail of the opening, such as the recess of the frame and the glazing, is retained on this cutout surface rather than on the wall.

\iffalse
\begin{figure}[hbht]
 \centering
 \includegraphics[width=1\textwidth]{Scene_dynamification_opening_cutout}
 \caption[Opening cutout from a wall mesh]{The detected opening polygon (left) is clipped out of the wall mesh triangle by triangle, producing a separate window or door surface (right) that keeps the original texture and plane.}
 \label{fig:8_opening_cutout}
\end{figure}
\fi


% ── 8.3.1.3 ──────────────────────────────────────────────────────────────────
\subsubsection{Dynamic Opening Construction}
\label{sec:8_opening_dynamic}

The extracted surface is still a flat patch coincident with the wall. To behave as a real, movable opening, it is given thickness and split into independently scalable parts (Figure~\ref{fig:8_opening_dynamic}). First, the outer perimeter of the cutout is extruded back toward the wall by a fixed depth, so the opening sits flush in the surface and gains the volume of a frame rather than reading as a decal. A boolean bounding box is then built around this volume and used to procedurally remove the part of the wall it intersects: rather than re-meshing the wall, the intersection is carved at render time through a screen-space stencil-buffer boolean~\cite{hecomi_screenspaceboolean_2018}, which keeps the operation interactive and free of geometric artefacts as the opening is moved or resized. Finally, the opening itself is separated into its framing and its panels or panes and 27-sliced~\cite{deftly_27_2021}, so that resizing stretches only the glazing and infill while the frame keeps its real thickness, matching how a physical window or door scales.

\iffalse
\begin{figure}[hbht]
 \centering
 \includegraphics[width=1\textwidth]{Scene_dynamification_opening_dynamic}
 \caption[Construction of a dynamic opening]{The flat cutout (left) is extruded to gain frame depth, carves the wall through a stencil-buffer boolean, and is 27-sliced into frame and panels (right) so it can be moved and scaled while the frame keeps its thickness.}
 \label{fig:8_opening_dynamic}
\end{figure}
\fi


% ── 8.3.1.4 ──────────────────────────────────────────────────────────────────
\subsubsection{Wall Re-Reconstruction}
\label{sec:8_wall_reinpaint}

Because windows and doors are now treated as separate, movable elements, the wall must be complete underneath them so that an opening can be repositioned or removed without exposing a gap. The wall surface behind each opening is reconstructed as a flat region flush with the surrounding wall: rather than reproducing the recessed geometry of the original opening, which now lives entirely on the cutout, the patch is kept planar so it matches the neighbouring wall and floor surfaces. Its texture is completed with the same per-plane strategy used in Chapter~\ref{ch:6_Scene_Completion}: the opening polygons are rasterised into an alpha mask on the wall texture, marking the opening interior as the region to fill, and the masked texture is completed by LaMa~\cite{suvorov_resolution-robust_2022} so that the wall continues plausibly behind the removed opening (Figure~\ref{fig:8_wall_reinpaint}). The wall mesh is closed across the hole using the same constrained triangulation as in Chapter~\ref{ch:6_Scene_Completion}, leaving a continuous structural surface with the opening represented as an independent overlay on top of it.

\iffalse
\begin{figure}[hbht]
 \centering
 \includegraphics[width=1\textwidth]{Scene_dynamification_wall_reinpaint}
 \caption[Wall re-reconstruction behind openings]{After the opening is removed (left), the wall surface behind it is flattened and re-inpainted (right) so that the structural shell remains complete when the opening is moved or removed.}
 \label{fig:8_wall_reinpaint}
\end{figure}
\fi


% ─────────────────────────────────────────────────────────────────────────────
\subsection{Object Dynamification}
\label{sec:8_object}

Depending on the completion method of Chapter~\ref{ch:7_Object_Completion}, each completed object arrives either as a textured mesh or as a mesh extracted from a completed \gls{SDF} with a texture applied. Object dynamification first brings both forms into a single deformable representation, the \gls{CSDF}, then turns it into a part-aware, selectively scalable field through convex decomposition, scaling-zone definition, joint interpolation, and, for large deformations, repetition of modular parts.


% ── 8.3.2.1 ──────────────────────────────────────────────────────────────────
\subsubsection{Part Segmentation}
\label{sec:8_part_segmentation}

Part segmentation is performed first on the completed mesh, because it produces an extra per-voxel attribute that is stored in the \gls{CSDF} alongside geometry and colour. The mesh is decomposed into approximately convex components using CoACD~\cite{wei_approximate_2022}, which recursively cuts the mesh with planes selected through a tree search that minimises a collision-aware concavity metric (Figure~\ref{fig:8_decomposition}). Adjacent components are merged where their union remains convex, reducing redundant cuts. Because the decomposition is purely geometric, it generalises to arbitrary reconstructed objects without relying on a class label, at the cost of parts that do not always coincide with semantic components~\cite{vermandere_adaptive_2026}.

\iffalse
\begin{figure}[hbht]
 \centering
 \includegraphics[width=1\textwidth]{Scene_dynamification_decomposition}
 \caption[Convex decomposition of a completed object]{The completed object mesh (left) and the approximately convex decomposition produced by CoACD~\cite{wei_approximate_2022} (right), with each part shown in a distinct colour.}
 \label{fig:8_decomposition}
\end{figure}
\fi

Each component is assigned a discrete part index $i$, which is later sampled onto the voxel grid so that every voxel knows which structural component it belongs to. These indices drive the part-aware deformation in Section~\ref{sec:8_interpolation}.


% ── 8.3.2.2 ──────────────────────────────────────────────────────────────────
\subsubsection{CSDF Creation}
\label{sec:8_csdf}

The completed mesh, its texture and its part labels are unified into a single \gls{CSDF} (Figure~\ref{fig:8_csdf}). A signed distance value $f(\mathbf{x})_\text{sdf}$ is sampled on a regular $128 \times 128 \times 128$ voxel grid, and the object texture is sampled onto the same grid so that each voxel also stores a colour $\mathbf{C} = (R, G, B)$. The part index $i$ from Section~\ref{sec:8_part_segmentation} is transferred to the grid by nearest-neighbour assignment, adding a third attribute per voxel. The grid is exported to a 3D texture sliced one pixel per voxel, where the RGB channels hold the colour and the alpha channel holds the signed distance value, with the part index carried in an auxiliary channel. This packed representation encodes geometry, appearance and structure in one volume and is the unit on which all subsequent deformation operates~\cite{vermandere_adaptive_2026}.

\iffalse
\begin{figure}[hbht]
 \centering
 \includegraphics[width=1\textwidth]{Scene_dynamification_csdf}
 \caption[CSDF creation]{The completed mesh, its texture and its part labels are sampled onto a common voxel grid and packed into a 3D texture: colour in the RGB channels, signed distance in the alpha channel, and the part index in an auxiliary channel.}
 \label{fig:8_csdf}
\end{figure}
\fi


% ── 8.3.2.3 ──────────────────────────────────────────────────────────────────
\subsubsection{Scaling-Zone Definition}
\label{sec:8_scaling_zone}

Selective scaling is controlled by a simple planar constraint. A scaling zone is delimited by two parallel planes along a chosen axis, a start plane and an end plane, and the user translates a third destination plane along that axis to set the target extent of the zone. Voxels before the start plane stay fixed, voxels beyond the end plane are rigidly translated to keep the rest of the object intact, and the voxels between the start and end planes are the ones that are deformed. This planar control is intuitive and behaves identically across objects, independent of their shape or class.


% ── 8.3.2.4 ──────────────────────────────────────────────────────────────────
\subsubsection{Joint Field Interpolation}
\label{sec:8_interpolation}

Within the scaling zone, a new voxel grid is allocated at the target extent and the field is resampled into it. The part index controls which voxels take part in the deformation: rather than scaling only the voxels strictly between the two planes, every part that is intersected by the planes is scaled as a whole. This grows the selection out to the natural boundaries of the sliced components, so the deformation follows the structure of the object instead of an arbitrary planar cut and produces a more natural-looking selection of voxels.

The signed distance and colour are continuous quantities and are interpolated linearly along the scaling axis. For two samples at positions $x_1$ and $x_2$, the value at an intermediate position $x$ is

\[
f(x)_\text{sdf} = f(x_1)_\text{sdf} + \frac{x - x_1}{x_2 - x_1}\,\big(f(x_2)_\text{sdf} - f(x_1)_\text{sdf}\big),
\]

\[
\mathbf{C}(x) = \mathbf{C}_1 + \frac{x - x_1}{x_2 - x_1}\,(\mathbf{C}_2 - \mathbf{C}_1).
\]

To limit computation and memory, interpolation is restricted to voxels near the surface, and the signed distance is locally re-evaluated after interpolation so that the zero level set remains consistent and free of discontinuities.

The part index $i$ is discrete and cannot be interpolated. Instead, the boundary between two adjacent parts is located along the scaling axis at the position where the index changes and the surface is closest, so that part transitions are preserved through the deformation rather than blurred:

\[
x_b = \argmin_{x} \big( |f(x)_\text{sdf}| \;\text{ where }\; i(x) \neq i(x + \delta x) \big),
\]

where $\delta x$ is a small sampling offset used to detect the label change. Figure~\ref{fig:8_interpolation} illustrates the joint interpolation of geometry, colour and part indices within a scaling zone.

\iffalse
\begin{figure}[hbht]
 \centering
 \includegraphics[width=1\textwidth]{Scene_dynamification_interpolation}
 \caption[Joint interpolation within a scaling zone]{Within the scaling zone, the signed distance field, colour and part indices are jointly resampled, producing a smooth and consistent deformation of both shape and appearance.}
 \label{fig:8_interpolation}
\end{figure}
\fi


% ── 8.3.2.5 ──────────────────────────────────────────────────────────────────
\subsubsection{Repetition-Based Scaling}
\label{sec:8_repetition}

For large deformations, stretching a region introduces visible distortion: shelves grow too tall, cushions become elongated, and repeated elements lose their proportions. For regions whose parts form a repeating pattern, the components are duplicated instead of stretched (Figure~\ref{fig:8_repetition}). Given an original repeatable length $L$ and a target length $L'$, the number of repetitions is

\[
n = \left\lfloor \frac{L'}{L} \right\rfloor,
\]

and the region is duplicated $n$ times along the scaling axis. The residual difference between $n L$ and $L'$ is absorbed by slightly scaling each repeated section using the interpolation scheme of Section~\ref{sec:8_interpolation}, and the boundary layers between adjacent repetitions are blended to avoid visible seams. This preserves the size of individual components while still reaching the target dimension~\cite{vermandere_adaptive_2026}.

\iffalse
\begin{figure}[hbht]
 \centering
 \includegraphics[width=1\textwidth]{Scene_dynamification_repetition}
 \caption[Repetition-based scaling]{For repeating structures, the selected region (left) is duplicated rather than stretched (right), preserving the proportions of each component during large deformations.}
 \label{fig:8_repetition}
\end{figure}
\fi

All of these deformations are applied in real time. The packed \gls{CSDF} is never edited destructively; it serves as the interpolation source from which the deformed field is sampled on the fly, and the result is continuously visualised as a surface extracted with Marching Cubes~\cite{lorensen_marching_1998}. The deformed grid is only written back to a 3D texture, optionally resampled to a fixed output resolution, when the object is exported.


% ─────────────────────────────────────────────────────────────────────────────
\subsection{Recombination into a Dynamic Scene}
\label{sec:8_recombination}

The dynamified openings and objects are placed back into the completed environment to form a single interactive scene in the Unity game engine. The structural shell and the openings are kept as explicit meshes, while each scalable object is visualised by extracting its surface from the deforming \gls{CSDF} with Marching Cubes~\cite{lorensen_marching_1998} every frame, so that geometry, colour and part structure update live as the field is edited. The scaled objects are positioned at the locations recovered during object isolation in Chapter~\ref{ch:5_Scene_Segmentation}, and the separated windows and doors remain as movable overlays on the re-reconstructed walls of Section~\ref{sec:8_wall_reinpaint}.

All three element classes are manipulated interactively from this interface. The primary structure is reshaped by dragging individual walls forward or backward, which extends the floor, ceiling and adjoining walls to follow. The secondary openings can be translated across their wall and resized through their 27-sliced frame, with the stencil-buffer boolean of Section~\ref{sec:8_opening_dynamic} re-carving the wall as they move. The tertiary objects can be repositioned and selectively scaled through the start, end and destination planes of their scaling zones. Every operation provides immediate visual feedback, turning the static reconstruction into an editable digital environment.


% ─────────────────────────────────────────────────────────────────────────────
\section{Experiments}
\label{sec:8_experiments}

The dynamification is evaluated in two parts. The scene-level results show complete dynamic environments assembled from the preceding pipeline stages, demonstrating openings that can be moved and objects that can be scaled in place. The object-level results compare the three scaling strategies on representative objects, focusing on the visual and structural quality of the deformation. The scene-level experiments use the completed environments of Chapter~\ref{ch:6_Scene_Completion} on the V-Scan~\cite{vermandere_synthetic_2026} and Matterport3D~\cite{chang_matterport3d_2017} datasets introduced in Chapter~\ref{ch:3_Data_Acquisition}, since these provide the reconstructed wall fragments with flattened textures required for opening detection. The object-level experiments use objects reconstructed and completed through Chapter~\ref{ch:7_Object_Completion}, together with a set of high-quality models with well-defined repeating structures to test the repetition strategy on clean geometry~\cite{vermandere_adaptive_2026}. All objects are represented as $128^3$ \glspl{CSDF}, and all scaling is performed and rendered interactively in the Unity environment of Section~\ref{sec:8_recombination}.


% ─────────────────────────────────────────────────────────────────────────────
\subsection{Scene-Level Results}
\label{sec:8_exp_scene}

Figure~\ref{fig:8_dynamic_scenes} shows complete dynamic environments produced by the full pipeline. The openings detected and cut from the walls behave as independent surfaces, and the re-reconstructed walls remain continuous when an opening is moved or removed. Combined with the in-place scalable objects and the draggable structural shell, the result is a scene whose primary, secondary and tertiary elements can all be manipulated after reconstruction.

\iffalse
\begin{figure}[hbht]
 \centering
 \includegraphics[width=1\textwidth]{Scene_dynamification_dynamic_scenes}
 \caption[Dynamic environments from the full pipeline]{Complete dynamic scenes assembled from the pipeline. Top: a V-Scan environment with separable openings and scalable objects. Bottom: a Matterport3D environment processed through the same stages.}
 \label{fig:8_dynamic_scenes}
\end{figure}
\fi

\todo{Add quantitative scene-level evaluation of opening detection on V-Scan, where the ground-truth window and door locations are known from the procedural room generator (Chapter~\ref{ch:3_Data_Acquisition}). Report detection precision, recall and IoU of the detected opening polygons against the ground-truth openings, broken down by environment type, mirroring the table format used in Chapters~\ref{ch:5_Scene_Segmentation} and~\ref{ch:6_Scene_Completion}. The current paper does not contain this experiment, so the data still needs to be generated.}


% ─────────────────────────────────────────────────────────────────────────────
\subsection{Object-Level Results}
\label{sec:8_exp_object}

Figure~\ref{fig:8_scaling_comparison} compares the three scaling strategies on representative objects. Global scaling uniformly resizes the whole object, so structural elements distort: chair legs become disproportionately wide and shelf spacings grow unnaturally. Selective scaling restricts the deformation to the scaling zone, preserving the proportions of components outside it, and performs well on simple parts such as legs or a tabletop. On objects with repeating structure, however, selective scaling stretches the repeated elements and produces visibly elongated, inconsistent geometry. Selective tiling resolves this by duplicating the repeating components instead of stretching them, so that shelves, cushions and similar elements keep their original size while the object reaches the target dimension.

\iffalse
\begin{figure}[hbht]
 \centering
 \includegraphics[width=1\textwidth]{Scene_dynamification_scaling_comparison}
 \caption[Comparison of scaling strategies]{Comparison of global scaling, selective scaling and selective tiling on reconstructed and high-quality objects. Global scaling distorts structural proportions, selective scaling stretches repeated parts, and selective tiling preserves component proportions by duplication.}
 \label{fig:8_scaling_comparison}
\end{figure}
\fi

\todo{Add the empirical object-level comparison promised in the thesis outline: quantify the three scaling methods rather than only showing them. The adaptive scaling paper~\cite{vermandere_adaptive_2026} is qualitative, so a metric needs to be chosen and computed, for example surface-distance or proportion-preservation of individual parts before and after scaling, or a user study on visual plausibility. State explicitly which metric is used and on how many objects per dataset.}


% ─────────────────────────────────────────────────────────────────────────────
\subsection{Discussion}
\label{sec:8_exp_discussion}

The experiments confirm the structure-aware ordering of the three strategies: global scaling is the least faithful, selective scaling improves proportionality outside the deformed region, and selective tiling is required once repeated structure is involved. Two limitations carry over from the underlying representation. The $128^3$ \gls{CSDF} resolution trades detail against memory, so fine geometry can be lost during interpolation or resampling, and the convex decomposition is geometry-based, so its parts do not always align with semantic components, which occasionally places scaling-zone boundaries at non-intuitive locations~\cite{vermandere_adaptive_2026}. At the scene level, opening detection inherits the planarity assumption of Chapter~\ref{ch:6_Scene_Completion} and the quality of the flattened wall texture: openings on curved or poorly inpainted walls are detected less reliably. These limitations are revisited in Chapter~\ref{ch:conclusion}.


% ─────────────────────────────────────────────────────────────────────────────
\section{Conclusion}
\label{sec:8_conclusion}

This chapter presented the dynamification stage that converts the static reconstructions of Chapters~\ref{ch:6_Scene_Completion} and~\ref{ch:7_Object_Completion} into a single interactive \gls{DRM}, acting on all three classes of building element. The primary structural shell is made draggable, the secondary windows and doors are detected in the flattened wall textures, separated as dynamic framed openings and re-inpainted behind, and the tertiary objects are converted into \glspl{CSDF}, decomposed into convex parts, and selectively scaled, either by jointly interpolating the signed distance, colour and part indices within a user-defined zone or, for repeating structures, by duplicating modular components, so that overall dimensions adapt while individual parts keep their proportions~\cite{vermandere_adaptive_2026}.

The qualitative results show that selective scaling preserves proportions far better than global scaling and that selective tiling is necessary to keep repeated structures coherent under large deformations, while the scene-level results demonstrate complete environments with separable openings and in-place scalable furniture rendered interactively in Unity~\cite{zhou_real-time_2008}. The remaining limitations, namely the voxel-resolution trade-off, the geometry-only part decomposition and the dependence on planar, well-textured walls for opening detection, define the directions for the future work discussed in Chapter~\ref{ch:conclusion}. With this stage, the pipeline closes the loop set out in Chapter~\ref{ch:introduction}: a partial scan is aligned, segmented, completed and finally made dynamic, yielding a complete and editable representation of the captured environment.


%%%%%%%%%%%%%%%%%%%%%%%%%%%%%%%%%%%%%%%%%%%%%%%%%%
% Keep the following \cleardoublepage at the end of this file,
% otherwise \includeonly includes empty pages.
\cleardoublepage